%\documentclass[a4paper,12pt]{amsart}
%\usepackage[dvipdfmx]{graphicx,xcolor}%graphics
%\usepackage{float}
%\usepackage{amsmath}
%\usepackage{verbatim}
%\usepackage{latexsym}
%\usepackage{amssymb}
%%\usepackage{emath}
%\usepackage{listings}
%\usepackage{ccicons}
%\usepackage{mflogo}
%%\usepackage{fancybox}%enclose by box
%\usepackage[dvipdfmx]{hyperref}
%\usepackage{verbatim}
%\hypersetup{% hyperrefオプションリスト
%   setpagesize=false,
%   bookmarksnumbered=true,%
%   bookmarksopen=true,%
%   colorlinks=true,%
%   linkcolor=red,
%   citecolor=blue
%} 
%
%\newtheorem*{Thm}{定理}
%\newtheorem*{Prop}{命題}
%
%\newtheorem*{exQ}{レポート}
%
%\theoremstyle{remark}
%\newtheorem*{ex}{例}
%\newtheorem*{Q}{演習}
%\newtheorem*{Quiz}{問題}
%
%\renewcommand{\figurename}{図}
%\renewcommand{\proofname}{証明}
%\renewcommand{\tablename}{表}
%\setlength{\belowcaptionskip}{0mm}
%
%\DeclareMathOperator{\Ker}{Ker}
%\DeclareMathOperator{\Image}{Im}
%
%\def\emptyline{\vspace{11pt}}
%\def\smallemptyline{\vspace{6pt}}
%\begin{document}
%
%\begin{flushright}
%   %2026/04/08   
%   2026/07/14
%   
%   九大・数理 大津幸男
%\end{flushright}
%
%
%\begin{center}
%   {\Large 情報科指導法II---11}
%   \href{https://creativecommons.org/licenses/by-nc-nd/4.0/}{\ccbyncnd}
%   \par{情報・社会・AI・健康}
%\end{center}

\chapter{情報・社会・AI・健康}
\section{情報とは}
 Data(データ)は数字や文字や図形などで表されたもの，
Information (情報) はデータに意味を持たせたもの，
Knowledge(知識)は情報を解釈して体系化したもの，
Wisdom(知恵)は知識を体得して人生に役立つものとしたものである．DIKW ピラミッドとはそれらを以下のように図式化したものである：
\begin{figure}[htb]
   \centering
   \includegraphics{~/Documents/GitHub/introInfo/ch6/images/figure.1}
   \caption{DIKW ピラミッド}\label{DIKW}
\end{figure}
 
 例えば生成AIにレポートを丸投げして内容を確認せず提出するのは，
 情報のみ得ることに対応する．出てきたレポートを自分なりに解釈し
 自分の言葉で喋れるようにするのは知識の習得にあたる．
 生成AIと壁打ちしながら自分に欠けていた視点からの考察を深め，
 自分だけでは書けなかったレベルのレポートをAIと共同で作成するのは
 知恵を得る使い方である．
 
%\section{情報社会とビジネス}
%近年ではビジネスや経営，マーケティングで情報技術を使うことが多くなった．
%
%CRM（Customer Relationship Management：顧客関係管理）

\section{情報社会と法律}
情報社会と情報セキュリティーに関係する日本の法律：
\begin{itemize}
   \item
      知的財産権

   \item
      著作権

   \item
      個人情報保護法
   \item      
      不正アクセス禁止法
   \item      
      サイバーセキュリティ基本法
      
   \item
      電気通信事業法
     \item    
      電子署名及び認証業務に関する法律
    \item     
      電波法
   \item      
      特定電子メールの送信の適正化等に関する法律
   \item      
      有線電気通信法
\end{itemize}


\section{AIの発達の歴史}
\begin{itemize}
   \item
      ロジックセオリスト：
      1955年から1956年にかけてアレン・ニューウェル，ハーバート・サイモン，
      J・C・ショーが開発したプログラム．
      人間の問題解決能力を真似するよう意図的に設計された世界初のプログラム
  
   \item 
      イライザ ELIZA：1964年ワイゼンバウムが開発した対話システム．
 
   \item  
      エキスパートシステム
   
%   チューリングテスト：
%   チューリングの考案した思考実験。機械と人間を別々の密室に入れ、第三者がそれぞれと交信して、機械と人間の判別がつかなければ、機械は人間と同じ思考能力をもつにいたったとする。
%   chatGPTにより不十分であることが示された．
   \item  
     neural network ニューラルネットワーク，
     machine learning 機械学習，
     deep learning  深層学習などのモデルが研究された．
 
    \item    
      ディープブルー： 1997年チェスの世界チャンピオンカスパロフに勝った
    \item    
      ワトソン： 2011年クイズ番組ジョパディで人間に勝利
   
%   ケンブリッジ・アナリティカ：
%   2016の大統領選でフェースブックのデータを使いトランプ大統領の勝利を導いた．アメリカで大問題となった．
    \item       
      AlphaGo 2016年碁で名人に勝利，  
      GAN =Generative Adversarial Network：敵対的生成ネットワークを使った
  
%   CNN たたみ込み型ニューラルネットワーク 
%   generative 生成的
%   adversarial 敵対的
%   generator 生成器
%   discriminator 識別器
%    RAG =retrieval Augmented genration
%   Claude 4 アンソロピック
%   Claude Mythos ミトス：

     \item
        LLM =large language model大規模言語モデルが発達し
        GPT-=Generative Pre-Trained Transformerが開発され，
         2023年それを対話に使うChatGPT が発表され
         世界に生成AIブームをもたらし，株もAIバブルの様相を示し，
         現代に至る．

%    OpenAI  Sam Altman CEO
%   ChatGPT 
%   GPT-3の次世代バージョン　
%   GPT-4 2023/3
   
\end{itemize}


\section{情報社会と心と体の健康}
最近，体やメンタルの不調を訴える人・社会的な問題を起こす人が増加している．

%
%{\Large 質問：}
%\begin{enumerate}
%	\item 
%	スマホを一日何時間使っていますか？
%	（iPhone ならScreen TimeでAndroidならDigital Wellbeing で確認できます）
%
%	\item 
%	いつも「ながらスマホ」していませんか？
%
%
%	\item 
%	猫背・ストレートネック・視力の低下・その他の目の不調・耳の不調・
%	頭痛・腰痛などありませんか？
%
%	\item 
%	運動はどれ位していますか?
%
%	\item 
%	一日何時間椅子に座っていますか？
%
%	\item 
%	一日何時間寝ていますか？
%
%	\item 
%	朝何時に起きますか？朝日を浴びていますか？
%
%	\item 
%	朝食は食べていますか？
%
%	\item 
%	コンビニで売っているような食べ物やファストフードをどれ位
%	食べていますか？
%\end{enumerate}
%
%\emptyline
%
\emptyline

\subsection{データ}
\begin{enumerate}
   \item 
      2023年のネット利用時間は高校生が6時間14分，中学生が4時間42分，
      小学生（10歳以上）が3時間46分だっただった
      \footnote{こども家庭庁 2023年度のインターネット利用環境実態調査}.
   
   \item 
      2024年の裸眼の視力が1.0未満の割合は高校生で71％，中学生で60％，
      小学生で36％だった
      \footnote{文部科学省 2024年度学校保健統計}．
   
   \item 
      2021年度中等度以上のうつ症状がある子どもの割合が高校生30％，
      中学生24％であると発表された
      \footnote{国立成育医療研究センター「コロナ×こどもアンケート」
      第4回調査報告}．
   
   \item 
      小中高校生の不登校の人数が2023年には過去最多の34万人を越えた
      \footnote{文部科学省 令和5年度児童生徒の問題行動・不登校等生徒指導上の
      諸課題に関する調査}．
   
   \item 
      2024年オーストラリアでは16歳未満のSNSを禁止する法律を定めた．

   \item 
      ネットカジノ依存
   \item 
      闇バイトの1年間の逮捕者の推移：
      年分トクリュウ全体の年間摘発者数うち「闇バイト」応募・関与者の人数2023年（令和5年）10,105人約3,925人（全体の約4割）2024年（令和6年）12,178人4,759人（全体の約4割）
\end{enumerate}
\emptyline


\subsection{対策の前提となる考え方}
\begin{enumerate}
  \item 
   科学的な（実験や統計や論理や数式で白黒付けられるような）
   反証可能(\cite{PK})な議論・
   対話可能な議論を行う．
   
   具体的対処の有効性は，体や心の状態・症状の改善度から評価する．
   特に，排便の状態や疲労感・睡眠不足感をモニターする．
%   \emptyline
   
   自分のことは自分では分かりづらい．
   例えば目の10cm 前でスマホを見ている人がいたら
   自分のスマホを目から離す行動をとるなど，
   他人を詳細に観察しそれが自分に当てはまらないか思い返してみるのは
   良いことである．
   
   \item 
   進化生物学：ヒトが農耕を始めてから未だ400世代程しかたっていないので，
   それ以来の生活の変化に我々の体は適応していない．
   ヒトが適応している1・２万年前の生理学的環境・状態と近い（
   逆に人工的な環境の）生活を送る必要がある(\cite{KL}, \cite{L})．
   
   \begin{enumerate}
      \item 
      ヒトは狩猟採集生活に適応しているが，運動しない生活，
      座り続ける生活，夜寝ない生活，
      長時間近くを見続ける生活，イヤホンを使い続ける生活
      には適応していない．
      
      \item 
      ヒトは狩猟採集生活の雑食の生活に適応しているが，
      多様な食物を摂らない生活，超加工食品には適応していない．
      
      \item 
      ヒトはスマホを使い続ける生活には適応していない
      (\cite{Han}, \cite{I}, \cite{Ko}, \cite{KS}, \cite{M2})．
   \end{enumerate}
	
   \item 
   動物は環境に合わせて行動を変えるように進化してきた．
   これを制御しているのが脳の報酬系
   \footnote{1954年オールズらにより最初に発見されたときは快楽中枢と
   呼ばれていたが，後に快楽ではなく，
   新しい行動や情報への報酬を与えていることが明かになり，
   報酬系と呼ばれるようになった．}
   から放出されるドーパミン\footnote{現在では脳内物質で
   説明されているが，古くからスキナーによる小動物を使った
   実験から現象としては知られていた．
   そのため，スマホをスキナー箱という研究者もいる．
   この心理学の分野を行動分析学\cite{M}という．}
   であり，
   ドキドキわくわくする感情を引き起こし，
   我々が新しい行動を行ったり新しい情報を得るのを促進するように働いている．
   ドーパミンはアドレナリンに代謝され攻撃性を増す効果もある．

   
   スマホ（のアルゴリズム）はドーパミンを継続的に大量に放出させ，
   使用者を依存症にするように作られており，
   それを作っている会社はそこから巨万の富を得ている．
   ただし，ドーパミン自体は善でも悪でもない．
   ドーパミンを上手くコントロールして，
   なりたい自分に近づいてほしい(\cite{LL})．
   
   \item
   我々のデジタルの行動データはブラウザーのCookie などに蓄積されていて，
   それが売買されたり，アルゴリズムがフィルターバブル
   \footnote{「インターネットの検索サイトが提供するアルゴリズムが，
   各ユーザーが見たくないような情報を遮断する機能」（フィルター）のせいで，
   まるで「泡」（バブル）の中に包まれたように，自分が見たい情報しか見えなく
   なること．(Wikipedia 2024/04/07, 18:34)} (\cite{P})
   ，SNSがエコーチェンバー現象
   \footnote{自分と似た意見や思想を持った人々の集まる空間
   （電子掲示板やSNSなど）内でコミュニケーションが繰り返され，
   自分の意見や思想が肯定されることによって，それらが世の中一般においても
   正しく，間違いないものであると信じ込んでしまう現象．
   (Wikipedia 2024/04/07, 18:35)}を引き起こしている．
   更に，最近ではこれらの現象に加えてアテンションエコノミーの拡大
   \footnote{情報の優劣よりも「人々の関心・注目」という希少性こそが
   経済的価値を持つようになり、それ自体が重要視・目的化・資源化・交換財化
   されるようになるという実態を指摘した概念．(Wikipedia 2025/03/21, 8:12)}
   や諸外国からの情報戦・認知戦により誤情報・偽情報が生み出され，
   それらが生成AI により著しく増幅されている．
   それらにさらし続けられることで，それらの中に所謂「真実」を見つけ
   \textcolor{red}{陰謀論に絡め取られる}人々が増えている(\cite{M1}, \cite{NKO})．
   陰謀論に惹かれることもドーパミンの放出と関連している．

   \item 
      それらによりトランプ大統領のようなポピュリズム政治家が世界中に台頭している．
    
\end{enumerate}
\emptyline

\subsection{デジタルデトックス}
\begin{enumerate}
%   \item 
%   スマホやデバイスの使用制限\footnote{\cite{PC}などを参考に．}
%   スマホやデバイスの使用禁止\footnote{\cite{PC}などを参考に．}---
%   これは一時的なもので，将来はスマホを自制して使えるようになることを
%   目標にする．
   \emptyline
   
   \item 
   生活改善：現在の生活はデジタルを除いても昔はなかった危険に囲まれている．
   もちろんスマホ＞ストレートネック＞肩こりのように直接の関連がある場合もある
   ので，どちらにしてもここから改善する必要がある．
   \begin{enumerate}
      \item 
      バランスのとれた（多様性の高い）質の高い食事を適量，よく噛んで食べる．
      \begin{enumerate}
         \item 
         野菜・果物・豆類・ナッツ・魚・肉・タマゴ・乳製品・
         発酵食品・ビタミン・ミネラルをバランスよく摂る．
         
         \item 
         コンビニで売っているような食品・ファストフードなどの
         超加工食品を食べない．これらにもドーパミンを放出させる働きがあり
         強い依存性を持っている．
         例：野菜ジュース・フルーツジュース・清涼飲料・カップ麺・
         ポテトチップス・揚げもの・加工肉・コンビニ弁当・サンドイッチ・
         ハンバーガー
         
         \emptyline				
         ただし，時たま友達や家族と一緒にこのようなものを食べるのは特に
         問題ない．逆に，他人に自分と同じ基準での食生活を求めるのは
         現実的でなく，場合によっては人間関係をこわすのでお奨めしない．
         
         \item 
         毎日朝食をとる．
         
         例：プチトマト・スプラウト・ゆで卵・納豆・豆腐・バナナ・ミカン・
         ヨーグルト・ゴマ・お酢・アマニ油・スープ・紅茶・コーヒー
         の中から何点かを適宜組み変えて摂る
         \footnote{基本パッケージ化されていて冷蔵庫を出して直ぐに食べられる
         もので(i)の要素を多数含むように選んでいる．}．
%         ご飯・味噌汁・焼きシャケ・卵焼き・キャベツの千切り・漬物
%         のような旅館の朝ご飯のようなもので続けられれば
%         それで構わないが，
%         長続きして(i)の要素を多数含むようなものを選んでほしい．
%         その意味では小型の食洗機の導入も検討に値する．}
      \end{enumerate}
   
      \item 
      バランスのとれた質の高い運動を適量，正しい姿勢で行う．
      \emptyline
   
      注意点：
      口呼吸の矯正，猫背・巻き肩・ストレートネックの矯正，
      肩甲骨の可動域の確保，
      反り腰の矯正，体の左右差の矯正，
      アキレス腱の柔軟性・アーチの保全・足指可動域の確保．
      \emptyline
   
      \begin{enumerate}
         \item 
         毎日1, 2時間ほど日光を浴び遠方を見ながらの
         ウォーキング・自転車・スポーツを行う．
         
         \item 
         長時間椅子に座り続けない．
         
         \item 
         階段の使用・片足立ち・発声練習・ストレッチ
         （専門家の指導が必要な場合は理学療法士にご相談ください）・
         足指のストレッチなどを隙間時間に行う．
      \end{enumerate}
      
      \item  
      質の高い睡眠と休養を適量，規則正しくとる．
      \begin{enumerate}
         \item 毎日8時間以上の規則正しい睡眠．
         \item 毎日朝日を浴びる．
      \end{enumerate}
      
      \item 
      バランスのとれた質の高い，社会活動を含む高度な知的活動を適量行う．
      
      \item 
      昔（殆ど）存在しなかった化学物質への被爆を避ける．
      
      %昔（殆ど）存在しなかった化学物質：
      砂糖・人工甘味料・合成洗剤(\cite{Uts})・大量のサラダ油(\cite{Ha})・
      トランス脂肪酸・PFAS・マイクロプラスチック・サプリ・抗生物質(\cite{Br})・
      薬など
      			
      \item 
%      フィルターバブル・エコーチェンバー・アテンションエコノミー・認知戦
%      について理解を深め，
      我々の世界認知がアルゴリズムや金儲けや敵国や
      生成AIから影響を受けていることを認識し，
      自分の認知が歪まされていないか常に検討することが必要である．
   \end{enumerate}
   	これらはできる限り続けることをお奨めする．
   \item
      デバイスとの付き合い方：
      \begin{enumerate}
            \item
            様々な通知機能は絶対外せないものを除いて停止．
            
            \item
            メール，SNS，ニュースの確認は一日数回に限定．週末は見ない
            と決めて知人に知らせるのもいい．
            
            \item
            集中しないといけないときはデジタルデバイスの電源を落とす．
            
%            \item
%            毎日1・2時間はスマホを見ずに，遠くの景色を見ながらの早歩きなどの
%            屋外活動を行う．日光がビタミンDを増やし近視の進行を止めることも
%            知られている．
%            ストレッチ・姿勢の矯正・片足立ち・階段の積極的使用などの運動なども
%            継続して行うことをお奨めする．
            
            \item
            スマホを含めデバイスを見るときはなるべく大きな外付けモニターを使い，
            表示フォントも大きなものを使う．
            モニターはなるべく遠くに置き，目に過度な負担をかけないようにする．
            iPhone では Setting /Screen Time/Screen Distance で画面との距離の警告を
            設定できる．
            また，
            イヤホンやヘッドホンではなく外付けのスピーカーを使い，
            他人にうるさくない音量に保つ．
            悪い姿勢を続けないように気を付ける．
            
            目も耳も体も（頭脳も）一生もの，大切に扱ってください．
      \item
         まめに\textcolor{red}{Cookie のリセット}などを行って
         フィルターバブルを防ぐことをお奨めする．
         詳しいやり方はブラウザー毎に検索するように．
      \item
         \textcolor{red}{プログラミングでデジタルデトックス}：
         プログラミングやデバッキング，論文やプレゼンの準備は
         アナログな作業．コードを書く前に，
         アナログでどれだけ綿密に計画を立てられるかで
         プロジェクトとその保守拡張の成功が決まる．
         データ構造とアルゴリズムをよく検討し，
         必要に応じて手で図を作成し，
         場合によっては数学的な性質を証明する必要がある．
         また，コーディングではパソコンでの作業が増えて
         スマートホンの小さな画面から離れられる．
         
         私の経験からのアドバイス：
         プログラミングに集中するとデジタルに距離が置けるようになる．
   \end{enumerate}

%   \emptyline
%   
%   \item 医療との向き合い方：
%   医療は必ずしも患者に良いことをやっているわけではない．
%   医療も商売の一つである．もちろん，
%   急激に体調が悪化した場合や大けがを負った場合直ぐに医療を受けることを勧める
%   が，医療とはほどほどの距離を保つことをお奨めする
%   (\cite{A},  \cite{Hi}, \cite{Me}, \cite{US})．
%   
%   \begin{itemize}
%      \item 商売で行われている治療も沢山ある．
%      		
%      \item 製薬会社のプロパガンダで作り上げられた病気もある．
%      	
%      \item 薬や治療が原因の不調も多い．
%      	
%      \item 体調が悪くないのに健康診断で病気を探すようなことは有害無益．
%   \end{itemize}
\end{enumerate}

%これをひと月続けて症状が改善するようなら，それを続けてください．
%改善しないようなら，やり方に問題なかったか反省して，やり方を変えて続けてください．それでもダメなら，既に生活改善では改善出来無いようなトラウマを抱えているか，全く別の原因がある可能性があります．
%その時こそ医師に相談することお奨めします．
%\emptyline
%
%ここに書いてあること全てを無批判に信じてはいけません．
%それではネットの情報を鵜呑みするのと変わりありません．
%以下に挙げた資料等を検討し必要ならそれ以外の文献も参考にし，
%心の底から納得したものがあればそれについてここに書いてあることを実践して下さい．



%\section{熱力学・生命・人間・社会・AI}
%\begin{itemize}
%   \item
%      エントロピーと熱力学の第二法則
%
%   \item
%      マックスウェル：マックスウェルの悪魔
%      
%   \item
%      シュレーディンガー：生命は負のエントロピーを食べて生きている．
%
%   \item
%      ショーペンハウエル：生命とは生きようとする盲目的な意志である．
%
%   \item
%      フォン・ノイマン：自己増殖オートマトン
%   \item
%      マルグリス：真核細胞の進化の内部共生説
%   \item
%      言葉・貨幣・神・法・文明の創成
%   \item
%      AI：人間と異質の知能      
%   \item
%      AIと人間との共存？
%\end{itemize}
%

%\section{授業で出来なかったこと}
%今回の授業で紹介できなかったこと：
%\emptyline
%
%基礎理論：
%\begin{itemize}
%   \item
%      サンプリング定理の証明
%   \item
%      計算可能性・再帰関数：チャーチの提言
%   \item
%      アルゴリズム
%   \item
%      計算量
%   \item
%      情報量
%\end{itemize}
%
%計算機：
%\begin{itemize}
%   \item
%      ノイマン型計算機のメカニズム
%   \item
%      オペレーティングシステム
%   \item
%      ユーザーインターフェース
%\end{itemize}
%
%プログラミング：
%\begin{itemize}
%   \item
%      Python 入門
%   \item
%      Shell script  入門
%   \item
%      Mail のプロトコル
%   \item
%      データベース, SQL入門
%   \item
%      プログラミング・プログラマーの思想と技法
%\end{itemize}
%
%
%統計と生成AI：
%\begin{itemize}
%   \item
%      データ分析：実践的技術
%   \item
%      統計理論
%   \item
%      機械学習
%   \item
%      LLM
%   \item
%      Fine turning
%   \item
%      生成AI への質問技術
%   \item
%      生成AI と著作権
%   \item
%      生成AI と兵器・戦争
%\end{itemize}
%
%これらについては，多くの本が書かれているので
%興味のある部分を自分で調べてください．

%\begin{thebibliography}{99}
%
%   \bibitem{Ba} 
%      ジェイミー バートレット, 
%      操られる民主主義: デジタル・テクノロジーはいかにして社会を破壊するか , 草思社
%
%   \bibitem{BF} 
%      ボズウェル, フィッチャー, リーダブルコード
%
%   \bibitem{B} 
%      ブレーザー, 失われてゆく我々の内なる細菌, みすず書店, 2015.
%   \bibitem{C}    
%      チャルディーニ：影響力の武器
%
%   \bibitem{H} 
%      アンデシュ・ハンセン, スマホ脳, 新潮新書, 2020.
%   \bibitem{H} 
%      ユバル・ノラ・ハラリ, サピエンス全史
%   
%   \bibitem{Ha} 
%      林 裕之, 「隠れ油」という大問題, 三五館, 2017.
%   
%   \bibitem{I} 
%      石井光太, ルポ スマホ育児が子どもを壊す, 新潮社, 2024.
%   
%   \bibitem{K} 
%      川本 晃司, スマホ失明, かんき出版, 2022.
%   
%   \bibitem{KL} 
%      ヴァイバー・クリガン=リード, 
%      サピエンス異変──新たな時代「人新世」の衝撃, 
%      飛鳥新社, 2018.
%   
%   \bibitem{KS}
%       川島 隆太, 榊 浩平, スマホはどこまで脳を壊すか, 朝日新書,  2023.
%   
%   \bibitem{L} 
%      ダニエル・リーバーマン, 人体六〇〇万年史 
%      上・下──科学が明かす進化・健康・疾病, 早川書房, 2015.
%   
%   \bibitem{LL} 
%      ダニエル・リーバーマン, マイケル・ロング , 
%      もっと! : 愛と創造，支配と進歩をもたらすドーパミンの最新脳科学, 
%   インターシフト (合同出版), 2020.
%
%   \bibitem{M} 
%   マロット, 行動分析学入門, 産業図書, 1998.
%
%   \bibitem{TH} 
%      達人プログラマー
%   
%   \bibitem{M1} 
%      物江 潤, デマ・陰謀論・カルト (新潮新書) , 新潮社, 2022.
%   
%   \bibitem{M2} 
%      物江 潤, デジタル教育という幻想 GIGAスクール構想の過ち, 
%      平凡社, 2023.
%   
%   \bibitem{NKO} 
%      長迫, 小谷, 大澤, SNS時代の戦略兵器 陰謀論 民主主義をむしばむ認知戦の脅威,
%    ウェッジ, 2025.
%   
%   \bibitem{KP} 
%      カール・ポッパー, 科学的発見の論理 上下, 恒星社厚生閣, 1971.
%   \bibitem{P} 
%      イーライ・パリサー, 閉じこもるインターネット——
%      グーグル・パーソナライズ・民主主義 , 早川書房, 2012.
%      改題：フィルターバブル──インターネットが隠していること
%      (ハヤカワ文庫NF), 早川書房, 2016.
%   \bibitem{PC}
%      キャサリン・プライス,  
%      スマホ断ち 30日でスマホ依存から抜け出す方法, 
%      角川新書, 2024.
%   
%   \bibitem{Ut} 
%   宇津木龍一, シャンプーをやめると髪が増える, KADOKAWA, 2013.
%\end{thebibliography}


%\end{document}

